% marc_aaai_appendix.tex — technical appendix for marc_aaai.tex (input after \bibliography).
% Under AAAI rules this is supplementary material; it does not count against the 7 content
% pages. Every table and number here is moved verbatim from marc.tex (port + cut, no rewrite).

\section{Per-Family Hybrid Battery and Cross-Family Transfer}
\label{app:hybrid}

The five implementation faults fixed before any comparison (the entrapment experiment): the denoiser never
received $x_t$; constants were absent from the graph tensors; $t$ did not condition variable
nodes; guidance magnitudes exploded; and a shape-handling fault in the one-variable case.
After the fixes, the solve rate on convex linear systems went from 0 to 1.000 in-distribution
and held out. A sixth, architectural fix is referenced from the Method section: per-round
LayerNorm washes out the magnitude of constraint constants, so a variable whose value is
pinned by its own constraint cannot be read off the deep features; a direct linear skip from
a variable's incident-factor constants to its output restores the path and cut the mean
absolute root error from 5.4 to 0.9.

Table~\ref{tab:hybrid} is the full mechanism ablation behind the hybrid-recipe study: four non-convex
families where deterministic descent is fully trapped (best-of-8, 60 held-out instances per
family; 95\% Wilson CIs shown on the random-init arm). The random-init control uses identical polish and budget with no
learning.

% provenance: RESULTS.md R3
\begin{table}[ht]
\centering
\caption{Mechanism ablation on trapped non-convex families (best-of-8, 60 held-out per
family; 95\% Wilson CIs on the random-init arm; columns are cold-start refinement, refinement with Langevin noise, random-init plus polish, multi-start Levenberg--Marquardt, and the learned hybrid). Values are those of the committed
\texttt{results/p\_hard/hard\_eval.json}. The multi-start Levenberg--Marquardt column is
reported for completeness: the strong classical solver saturates all four families, which is
why the learned-versus-random comparison, not the learned-versus-Langevin one, is the
informative contrast.}
\label{tab:hybrid}
\footnotesize
\setlength{\tabcolsep}{4pt}
\begin{tabular}{lccccccl}
\toprule
Family & cold & +Langevin & random-init & LM & learned & learned $>$ random? \\
\midrule
BilinearSystem  & 0.000 & 0.300 & 0.550 [0.42, 0.67] & 1.000 & 0.550 & tie ($p{=}0.50$) \\
BilinearProduct & 0.000 & 0.100 & 0.683 [0.56, 0.79] & 1.000 & 0.683 & tie ($p{=}0.50$) \\
QuadraticSystem & 0.000 & 0.300 & 0.683 [0.56, 0.79] & 1.000 & 0.683 & tie ($p{=}0.50$) \\
CircleLine      & 0.000 & 0.033 & 0.200 [0.12, 0.32] & 1.000 & 0.000 & no (random wins) \\
\bottomrule
\end{tabular}
\end{table}

\label{app:transfer}
Cross-family transfer is partial. We train on three non-convex families and test the hybrid
on the held-out fourth (best-of-8, 60 test instances per family, 200 epochs);
Table~\ref{tab:transfer} reports the outcome. On two of four held-out families the
cross-trained model reaches 0.683, significantly above the refine+Langevin baseline
($p < 10^{-4}$): the model solves a family it never trained on. On the other two it solves
nothing. The BilinearSystem failure is the informative one, because that family is solvable
in-distribution (0.550 in Table~\ref{tab:hybrid}). A smoke run trained on only the two
similar families \{Product, Quadratic\} recovered it at 0.70, whereas adding the pathological
CircleLine family to the training mix collapsed it to 0.00. Dissimilar training families can
degrade transfer. We read this as transfer across related non-convex structure, not
as a universal solver.

% provenance: RESULTS.md R4
\begin{table}[ht]
\centering
\caption{Leave-one-family-out transfer (best-of-8, 60 test per family, 200 epochs). $p$ from
a $z$-test of learned (cross) $>$ refine+Langevin.}
\label{tab:transfer}
\small
\begin{tabular}{lllll}
\toprule
Held-out family & Trained on & refine+Langevin & learned (cross) & $p$ \\
\midrule
BilinearProduct & Sys, Quad, Circle  & 0.100 [0.05, 0.20] & 0.683 [0.56, 0.79] & $<10^{-4}$ \\
QuadraticSystem & Sys, Prod, Circle  & 0.300 [0.20, 0.43] & 0.683 [0.56, 0.79] & $<10^{-4}$ \\
BilinearSystem  & Prod, Quad, Circle & 0.300 [0.20, 0.43] & 0.000 & fails \\
CircleLine      & Sys, Prod, Quad    & 0.033 [0.01, 0.11] & 0.000 & fails \\
\bottomrule
\end{tabular}
\end{table}

\section{Dimension Scaling: Full Battery}
\label{app:scaling}

Table~\ref{tab:scaling} is the full five-arm battery behind the dimension-scaling study; the main text and
Figure 1 of the main paper carry the two arms that decide the claim (random restart and learned).

% provenance: RESULTS.md R15 (methodology unified-v2, results/p_scaling/scaling.json)
\begin{table}[ht]
\centering
\caption{Dimension scaling on bundled traps with per-instance-varying roots in $\pm[3,8]$
(best-of-8, unified-v2 protocol: one shared polish and one checker for every arm, $N{=}40$
per cell; 95\% Wilson half-widths are at most $0.15$). Bold marks the best method per
row. A three-seed rerun of the same protocol ($N{=}120$ per cell,
scaling\_3seed.json) reads $0.983 \pm 0.014$ at $n{=}6$ and $0.658 \pm 0.484$ at
$n{=}4$: single cells are seed-noisy, the crossover is not.}
\label{tab:scaling}
\begin{tabular}{cccccc}
\toprule
$n$ & deterministic & Langevin & mean-prior & random restart & learned \\
\midrule
1 & 0.000 & 0.775 & 0.000 & \textbf{1.000} & 0.950 \\
2 & 0.000 & 0.225 & 0.000 & 0.725 & \textbf{0.950} \\
3 & 0.000 & 0.000 & 0.000 & 0.075 & \textbf{0.975} \\
4 & 0.000 & 0.000 & 0.000 & 0.000 & \textbf{0.925} \\
6 & 0.000 & 0.000 & 0.000 & 0.000 & \textbf{0.250} \\
\bottomrule
\end{tabular}
\end{table}

\section{The Factorization Dichotomy: Regime Map and Table}
\label{app:law}

The regime map (Figure 2 of the main paper) places every measured family on the
two axes of the law; Table~\ref{tab:law} tabulates the dichotomy plotted in
Figure 1 of the main paper (600 fresh instances per $n$, Wilson CIs, $K = 8$). Expected restarts
$1/q(n)$ are listed over the tested dimensions in increasing order.

Slope provenance for the regime map: the separable and coupled abscissae are the fits of
Table~\ref{tab:law}, the geometry abscissa is the fit of Figure 1 of the main paper, and the
three crossover-replication families are placed at slopes inverted from their best-of-8
Levenberg--Marquardt arm through the best-of-$K$ identity of the main paper. The real-systems suite is plotted
at a nominal abscissa: it has no learned arm and no separately measured $q(n)$.


% provenance: RESULTS.md R9 + paper/notes/crossover_law.md §5
\begin{table}[ht]
\centering
\caption{The factorization dichotomy, measured.}
\label{tab:law}
\small
\begin{tabular}{lll}
\toprule
Quantity & Independent (separable) & Coupled \\
\midrule
$\log q(n)$ slope & $-1.03$ & $-0.13$ \\
fit $R^2$ & 0.98 & 0.96 \\
measured $v = q(1)$ & 0.27 & --- \\
$E[\text{starts}] = 1/q(n)$ & $3.7 \to 13 \to 32 \to 150 \to 600$ & $2.0 \to 2.5 \to 2.9 \to 3.8 \to 4.3$ \\
$P_{\text{random}}$ MAE, parameter-free & 0.012 & law N/A (not separable) \\
\bottomrule
\end{tabular}
\end{table}

\section{External Anchor: Per-Class Construction Repair}
\label{app:anchor}

Table~\ref{tab:anchor} is the per-class breakdown behind the external-anchor paragraph.
Both restart controls rerun the same classical reference solver on the unrepaired system:
\texttt{restart} at the reference solver's own budget, \texttt{restart}$\times|V|$ at the
full enumeration budget over the construction vocabulary. The two classes that never fail
are reported as negatives rather than averaged in as zeros; the own-budget column shows how
far short a plain retry falls even before the budget is matched.

% provenance: RESULTS.md R30; results/p_real_repair/real_repair_multiseed.json (aggregate)
\begin{table}[ht]
\centering
\caption{Derived-construction repair on hardened variants of named real systems: three
independent seed bases, $N{=}200$ instances per class per seed, mean $\pm$ population SD.
``Failures'' is the two-stream failure rate of the classical reference and $N$ the pooled
failure population it defines; the construction is chosen on a disjoint half of that pool
and graded on the held-out half.}
\label{tab:anchor}
\small
\begin{tabular}{lrrccc}
\toprule
Class & Failures & $N$ & construction (held out) & restart & restart$\times|V|$ \\
\midrule
Trilateration, far side & $0.848 \pm 0.020$ & 509 & $\mathbf{1.000 \pm 0.000}$
  & $0.079 \pm 0.017$ & $0.433 \pm 0.049$ \\
Conic--line, ghost root & $0.263 \pm 0.015$ & 158 & $\mathbf{1.000 \pm 0.000}$
  & $0.026 \pm 0.010$ & $0.114 \pm 0.011$ \\
3R inverse kinematics & $0.000$ & 0 & --- & --- & --- \\
Circles, far apart & $0.000$ & 0 & --- & --- & --- \\
\bottomrule
\end{tabular}
\end{table}

\section{Repair: Menu-Size Scaling and Cost Accounting}
\label{app:kscaling}

The $K{=}4$ checkpoint transfers to larger linear menus without retraining, but the accuracy
advantage shrinks with $K$ and is gone by $K{=}16$ ($0.300/0.187/0.113$ against random
$0.227/0.120/0.107$ at $K{=}4/8/16$, $N{=}300/150/150$); what survives at large $K$ is only
the cost gap, as the enumeration it displaces grows to $9.05$ solver calls per instance
(Figure~\ref{fig:repair}). Directly training at $K{=}16$ performs at chance, so the
transferred checkpoint is selected and both negatives are kept in the record.

% provenance: paper/tex/figures/fig_repair.pdf (scripts/plot_repair.py), RESULTS.md R10
\begin{figure}[t]
\centering
\includegraphics[width=0.85\textwidth]{figures/fig_repair.pdf}
\caption{Left: Table 3 of the main paper drawn, the ranker against its
candidate-only and random controls, with $K{=}4$ chance dotted. Right: menu-size scaling.
The $K{=}4$ checkpoint evaluated zero-shot at larger menus retains an accuracy edge at
$K{=}8$ that closes by $K{=}16$, while the blind enumeration it displaces grows from $2.5$
to $9.1$ solver calls per instance.}
\label{fig:repair}
\end{figure}

\section{MATH-500 Scope Measurement}
\label{app:math}

The main text states that the template formalizer covers none of a 48-problem MATH-500
sample \citep{hendrycks2021math}. The breakdown: roughly $20\%$ of the sample is
constraint-shaped and in principle in scope, $31\%$ is CAS computation rather than
constraint solving, and $48\%$ is reasoning or proof that the factor-graph encoding does
not express. Coverage is $0/48$ because the constraint-shaped fraction still requires
autoformalization from problem text, which the template formalizer does not perform. The
binding constraint on applying this substrate to natural problem statements is therefore
autoformalization, not the solver.
% provenance: RESULTS.md R6 / PROVENANCE R14 (scripts/run_math_coverage.py)
